Au-delà des résultats obtenus, ce travail ouvre sur plusieurs axes d'amélioration, tant sur le plan méthodologique qu'algorithmique.

\subsection{Correction de l'implémentation SAC}

La problématique de ce travail interroge notamment le choix et la robustesse des algorithmes face à un environnement à récompense rare (Section~\ref{sec:problematic}). Or, la comparaison menée entre PPO et SAC en Section~\ref{sec:sac_comparison} n'a pas permis de trancher cette question de façon définitive pour SAC : l'interruption de son entraînement dès le Palier~I 
s'explique en partie par des limites d'implémentation de \textit{rlgym-sac} plutôt que par une caractéristique intrinsèque de l'algorithme -- absence de gradient clipping, standardisation des retours non appliquée à l'entraînement, cible d'entropie non paramétrable (Section~\ref{sec:limites-sac}).

Corriger ces limites constituerait donc un préalable nécessaire avant de pouvoir statuer réellement sur la capacité de SAC à surmonter la rareté du signal de récompense dans cet environnement, plutôt que de conclure prématurément à son inadéquation. Une telle correction permettrait notamment de déterminer si la stagnation observée relève uniquement des défauts logiciels identifiés, ou si elle révèle une difficulté plus structurelle des méthodes off-policy à entropie maximale à exploiter un signal de récompense rare dans un espace d'actions continu de grande dimension. Ce travail représente cependant un investissement d'ingénierie logicielle non négligeable, hors du périmètre initial de ce mémoire, dont l'objectif portait sur la comparaison empirique des algorithmes plutôt que sur le développement ou la maintenance de leur implémentation.

\subsection{Apprentissage par imitation à partir de vidéos}


L'ensemble des stratégies mises en \oe uvre dans ce travail pour répondre au problème du \textit{sparse reward} -- \textit{reward shaping} et \textit{curriculum learning} (Sections~\ref{sec:reward_shaping}, \ref{sec:curriculum}) -- agissent sur la fonction de récompense ou sur la progression de l'entraînement, mais laissent inchangée la façon dont l'agent explore initialement l'environnement : celui-ci part toujours d'un comportement aléatoire, et doit découvrir par essais et erreurs les comportements élémentaires (contact avec la balle, déplacement cohérent) avant même de pouvoir exploiter les signaux de récompense les plus denses. Une piste alternative, complémentaire à celles explorées dans ce mémoire, consisterait à initialiser l'agent avec un comportement déjà informé, plutôt que de le laisser découvrir ces fondamentaux depuis une politique aléatoire.

\textcite{bakerVideoPreTrainingVPT2022} démontrent la viabilité de cette approche sur Minecraft, un environnement partageant plusieurs caractéristiques avec Rocket League : un espace d'actions complexe, un objectif final éloigné des actions élémentaires, et une rareté structurelle du signal de récompense\footnote{Dans le papier, l'objectif est de fabriquer une pioche en diamant, une tâche nécessitant une longue séquence de sous-objectifs et requérant, pour un joueur humain compétent, une durée médiane supérieure à 20 minutes (environ 24\,000 actions).}. Leur méthode, le \textit{Video PreTraining} (VPT), consiste à entraîner un modèle de dynamique inverse sur un faible volume de données étiquetées, afin d'annoter automatiquement un grand volume de vidéos non annotées disponibles en ligne. Les séquences d'observations et d'actions ainsi reconstituées permettent d'entraîner un modèle comportemental par apprentissage supervisé, lequel peut ensuite être affiné par apprentissage par renforcement sur la tâche cible.

Cette approche a d'ailleurs déjà été adaptée à Rocket League par la communauté RLGym \parencite{rolv-arildRolvArildReplaypretraining2026} : un modèle de dynamique inverse y est entraîné à partir des entrées et sorties de bots existants tels que Nexto, puis appliqué à des fichiers de replay de joueurs humains pour en reconstituer les actions sous-jacentes, produisant ainsi un jeu de données exploitable pour un préentraînement par imitation avant affinement par RL.

Une telle approche pourrait accélérer l'acquisition des comportements fondamentaux du Palier~I (Section~\ref{sec:palier1}), potentiellement au bénéfice de SAC, dont l'échec à dépasser ce palier est en partie attribué à la difficulté d'explorer efficacement un espace d'actions continu de grande dimension (Section~\ref{sec:sac_comparison}). Elle comporte cependant un risque inverse à celui du reward shaping : plutôt que d'introduire des failles par une spécification imparfaite de la récompense (\textit{reward hacking}, Section~\ref{sec:reward_hacking}), l'imitation risquerait d'importer directement les inefficacités propres au comportement humain plutôt que de les corriger. L'analyse qualitative menée en Section~\ref{sec:ppo_results} illustre précisément ce risque : l'agent entraîné par RL pur a appris à couper le boost dès que l'accélération supplémentaire ne produit plus d'effet utile une fois la vitesse maximale atteinte (Figure~\ref{fig:boost_while_supersonic}), un comportement plus économe que celui des joueurs humains testés, qui tendent à maintenir le boost activé au-delà de ce seuil. Un agent préentraîné par imitation sur des replays humains risquerait ainsi de reproduire cette même inefficacité plutôt que de la dépasser, illustrant la nécessité de conserver une phase d'affinement par renforcement suffisamment longue pour corriger de tels biais hérités de la démonstration humaine.

\subsection{Entraînement multi-agent}

Ce travail s'est volontairement restreint au mode 1v1 (Section~\ref{sec:action_space}), afin d'isoler les défis de l'apprentissage individuel sans introduire la complexité supplémentaire de la coordination entre coéquipiers. Deux limites découlent directement de ce choix, chacune ouvrant sur une piste d'amélioration distincte.

\subsubsection{Généralisation à plusieurs modes de jeu}

Comme discuté en Section~\ref{sec:ppo_results}, la comparaison contre Necto 
\parencite{rolv-arildRolvArildNecto2026} appelle une nuance : cet agent a été entraîné simultanément sur les trois modes compétitifs (1v1, 2v2 et 3v3), tandis que l'agent produit dans ce travail est spécialisé exclusivement en 1v1. Si cette spécialisation constitue un avantage ponctuel dans le cadre de cette comparaison précise, elle limite la portée pratique de l'agent, qui ne pourrait affronter un adversaire humain en configuration 2v2 ou 3v3 sans réentraînement 
complet. Une extension naturelle de ce travail consisterait à étendre le curriculum développé en Section~\ref{sec:curriculum} à un entraînement conjoint sur les trois modes de jeu, à la manière de Necto \parencite{rolv-arildRolvArildNecto2026} ou d'OpenAI Five \parencite{openaiDota2Large2019}, ce dernier ayant démontré la faisabilité d'une coordination multi-agent complexe (5 contre 5) à grande échelle par self-play. 
Une telle extension poserait cependant la question de la conception des récompenses de coordination inter-agents (par exemple, encourager un partage cohérent des rôles offensif/défensif entre coéquipiers), qui n'a pas été abordée dans ce travail centré sur le 1v1.

\subsubsection{Diversité des partenaires d'entraînement}

Par ailleurs, l'agent produit dans ce travail a été entraîné exclusivement par \textit{self-play} contre une copie de sa propre politique (Section~\ref{sec:reward_hacking}). Si cette approche a permis d'éviter les stratégies collaboratives dégénérées identifiées en Section~\ref{sec:reward_hacking} 
une fois les récompenses reformulées en \textit{zero-sum}, elle expose l'agent au risque inverse : converger vers un ensemble restreint de stratégies bien adaptées à son propre style de jeu, mais potentiellement moins robustes face à la diversité des styles rencontrés chez des adversaires humains ou d'autres agents. \textcite{vinyalsGrandmasterLevelStarCraft2019} illustrent une 
alternative à ce problème avec le \textit{league training} : plutôt qu'un unique adversaire-miroir, AlphaStar est entraîné contre une population diversifiée d'agents en compétition permanente, incluant des agents dits \textit{exploiteurs} spécifiquement chargés d'identifier et de contrer les faiblesses des stratégies dominantes. Une adaptation de ce principe à Rocket League -- par 
exemple en confrontant périodiquement l'agent à des versions antérieures de lui-même ou à des agents tiers comme Element et Necto (Section~\ref{sec:ppo_results}) plutôt qu'à sa seule version courante -- pourrait renforcer la robustesse et la diversité stratégique de l'agent produit.